\documentclass[border=10pt]{standalone}
\usepackage{tikz}
\usepackage{amsmath}
\usetikzlibrary{arrows.meta, positioning, shapes.geometric, calc}

\begin{document}

\definecolor{pcolor}{RGB}{0, 120, 200}    % Blue - P
\definecolor{icolor}{RGB}{40, 160, 40}    % Green - I
\definecolor{dcolor}{RGB}{220, 50, 30}    % Red - D
\definecolor{arrowcolor}{RGB}{80, 80, 80}
\definecolor{sumcolor}{RGB}{120, 80, 150} % Purple - sum
\begin{tikzpicture}[
    block/.style={
      rectangle, draw=#1, line width=1.5pt,
      fill=#1!8, rounded corners=3pt,
      minimum width=38mm, minimum height=13mm,
      font=\bfseries, text=#1, align=center
    },
    % Summing-junction circle diameter is ~40% of block height (Block
    % Diagram Style convention), not 77% as before (10mm / 13mm block).
    % 加え合わせ点の円の直径はブロック高さの約4割（ブロック線図の流儀）。
    % 以前は 10mm / ブロック高さ13mm = 77% だった。
    sum/.style={
      circle, draw=sumcolor, line width=1.5pt,
      fill=sumcolor!8, minimum size=5.5mm,
      font=\normalsize\bfseries, text=sumcolor
    },
    arrow/.style={-{Stealth[length=3mm]}, line width=1.2pt, arrowcolor},
  ]

  % Input
  \node[font=\bfseries, arrowcolor] (input) at (-1, 0) {$e(t)$};

  % Branch point
  \coordinate (branch) at (1.5, 0);
  \fill[arrowcolor] (branch) circle (2.5pt);

  % P block (top)
  \node[block=pcolor] (P) at (5.5, 2) {P: $K_p \cdot e$};

  % I block (center)
  \node[block=icolor] (I) at (5.5, 0) {I: $\dfrac{K_p}{T_i} \displaystyle\int e\,dt$};

  % D block (bottom) -- Incomplete derivative filter
  \node[block=dcolor] (D) at (5.5, -2) {D: $K_p \cdot \dfrac{T_d s}{\eta T_d s + 1}$};

  % Sum
  \node[sum] (sum) at (10, 0) {$\Sigma$};

  % Output
  \node[font=\bfseries, arrowcolor] (output) at (12.5, 0) {$u(t)$};

  % Input arrow to branch
  \draw[arrow] (input) -- (branch);

  % Branch to P/I/D
  \draw[arrow, pcolor] (branch) |- (P);
  \draw[arrow, icolor] (branch) -- (I);
  \draw[arrow, dcolor] (branch) |- (D);

  % P/I/D to sum
  \draw[arrow, pcolor] (P) -| (sum);
  \draw[arrow, icolor] (I) -- (sum);
  \draw[arrow, dcolor] (D) -| (sum);

  % Sum to output
  \draw[arrow] (sum) -- (output);

  % Plus signs at sum.
  % Fix: the I-branch "+" was at 180 deg, exactly on the horizontal I-sum
  % arrow's own axis, so it was hidden under the arrowhead. Moved to
  % 160 deg (off-axis, slightly above) with a larger radius for clearance.
  % 修正: I分岐の「+」は180度（I→Sumの水平矢印の軸線上）にあり、矢頭の
  % 下に隠れていた。矢印の軸から外れる160度（軸よりやや上）へ移動し、
  % 半径も少し広げて余裕を持たせた。
  \node[font=\scriptsize, pcolor] at ($(sum)+(120:0.7)$) {$+$};
  \node[font=\scriptsize, icolor] at ($(sum)+(160:0.85)$) {$+$};
  \node[font=\scriptsize, dcolor] at ($(sum)+(240:0.7)$) {$+$};

  % Labels
  \node[font=\small, pcolor, anchor=south west] at (7.5, 2) {Proportional};
  \node[font=\small, icolor, anchor=south west] at (7.5, 0) {Integral};
  \node[font=\small, dcolor, anchor=north west] at (7.5, -2) {Incomplete Derivative};

  % Transfer function (below diagram)
  \node[font=\normalsize, arrowcolor, anchor=north] at (5.5, -3.8) {%
    $C(s) = K_p\!\left(1 + \dfrac{1}{T_i s} + \dfrac{T_d s}{\eta T_d s + 1}\right)$};

\end{tikzpicture}
\end{document}
